\section{Experimental Results}
\subsection{Experimental Settings}

\mypara{Datasets.}
We validated the effectiveness of \sysname\ on the two real-world autonomous driving datasets, the KITTI-360~\cite{Liao2022PAMI} and nuScenes \cite{nuscenes2019} datasets.
The KITTI-360 dataset features a Velodyne-64E LiDAR with 64 beams and 360-degree horizontal FOV. Following existing works~\cite{nakashima2024lidar, wu2024text2lidar}, we designated sequences 0 and 1 as the test set (26367 frames) and utilized the remaining sequences as training set (50348 frames). 
The nuScenes dataset provides 32-beam LiDAR measurements with 360-dergree horizontal FOV  at 20$\si{Hz}$ and keyframe annotations at 2$\si{Hz}$. 
Following common practice \cite{gao2023magicdrive}, we expanded the nuScenes bounding box labels to 12$\si{Hz}$ and adhered to the official training and validation dataset splits (700 training scenes and 150 validation scenes), resulting in 165280 training samples and 35364 validation samples.



\subsection{Unconditional LiDAR Scene Generation}
\label{sec_single_frame}

We compare \sysname\ against SOTA LiDAR scene generation methods,
LiDARGen~\cite{zyrianov2022learning}, R2DM~\cite{nakashima2024lidar}, Text2LiDAR\cite{wu2024text2lidar}, and RangeLDM~\cite{hu2025rangeldm}. 
\cref{tabel1} illustrates the unconditional generation results on the KITTI-360 test set and the nuScenes val split. 
It can be seen that \sysname\ achieves comparable generation quality than RangeLDM on the KITTI-360 dataset, with superior FRD \cite{wu2024text2lidar} and slightly inferior MMD \cite{nakashima2024lidar} and JSD \cite{nakashima2024lidar}. However, \sysname\ surpasses RangeLDM on the nuScenes dataset in all three metrics, FRD, MMD, and JSD. Overall, these results demonstrates superiority of \sysname\ among all compared methods.


\subsection{Conditional LiDAR Scene Generation}
\label{sec_conditional_generation}

\begin{table}[ht]
\centering
%\refstepcounter{table}
\scalebox{0.95}{
\begin{tabular}{lcc|ccc} 
\toprule
  & \begin{tabular}[c]{@{}c@{}}Road \\sketch\end{tabular} & \begin{tabular}[c]{@{}c@{}}Scene \\caption\end{tabular} & FRD$\downarrow$    & MMD$^1$$\downarrow$  & JSD$\downarrow$    \\ 
\midrule
Config A &                                                       &                                                         & 210.55 & 1.84 & 0.045  \\
Config B & \checkmark                                             &                                                         & 178.24 & 0.85 & 0.041  \\
Config C & \checkmark                                             & \checkmark                                               & \textbf{175.17} & \textbf{0.85} & \textbf{0.039}  \\
\bottomrule
\multicolumn{5}{l}{\footnotesize 1: MMD has been multiplied by \({10^{4}}.\)} \\
\end{tabular}}
\caption{
Conditional generation results of \sysname\ with different conditions on the nuScenes val split.}
\label{tab_ablation}
\end{table}


\begin{figure}[t]
    \begin{center}
       \includegraphics[width=\linewidth]{figures/ablation.pdf}
    \end{center}
       \caption{LiDAR scenes generated by \sysname\ with different conditions on the nuScenes val split.}
       \label{fig_ablation}
\end{figure}


We examine the impacts of road sketch and scene caption conditions by training \modelname\ with subsets of the multimodal conditions.
\cref{tab_ablation} presents results of \sysname\ with different conditions on the nuScenes val split.
Obviously, the influence of road sketch conditioning is particularly significant, yielding a 15.34\% improvement in FRD (from 210.55 to 178.24) and a 53.80\% enhancement in MMD (from 1.84 to 0.85) compared to unconditional generation.
This is attributed to the precise information about the road structure and the foreground objects embedded within the road sketches. Scene captions contribute to further enhancing the realism of the generated point clouds by providing complementary contextual information.


\cref{fig_ablation} shows the qualitative comparison of the LiDAR scenes generated by \sysname\ with different conditions on the nuScenes val split. 
Without any conditioning, the generated LiDAR scene lacks specific object structures (as shown in \cref{fig_ablation}b). 
When conditioned solely on road sketches (as shown in \cref{fig_ablation}c), \sysname\ successfully controls object locations in the generated scene, though the background area (\textcolor{red}{red circle 1}) is incomplete. 
Upon incorporating scene caption conditioning (shown in \cref{fig_ablation}d), both foreground objects (\textcolor{red}{red circle 2}) and background elements (\textcolor{red}{red circle 1}) exhibit more complete structures, demonstrating the complementary benefits of this additional conditioning. This improvement stems from the comprehensive nature of scene captions, which encapsulate both foreground and background contextual information.


\begin{figure}[ht]
    \begin{center}
       \includegraphics[width=\linewidth]{figures/fine_grained_control.pdf}
    \end{center}
       \caption{Visualization of fine-grained scene manipulation of \sysname\ on the nuScenes val split. Top row: road sketches. Bottom row: LiDAR scenes. a) GT scene, b) Generated scene with multimodal conditions, c) Generated scene with \textit{edited} multimodal conditions.}
       \label{fig_fine_grained_control}
\end{figure}

Furthermore, \cref{fig_fine_grained_control} demonstrates an example of fine-grained manipulation of road elements in a generated LiDAR scene. In \cref{fig_fine_grained_control}b, we re-simulated the real scene with same multimodal conditions as those of GT scene, where we can see that the generated scene closely align with GT scene (e.g., see buildings). In \cref{fig_fine_grained_control}c, we inserted an vehicle and changed the background. This manipulation is achieved through edited multimodal conditions comprising 1) addition of a 3D box to the road sketch, 2) extraction of he corresponding object prior, and 3) modification of the scene caption from \{\textit{... building ...}\} to \{\textit{... grassy area with trees...}\}. The resulting scene successfully incorporates these manipulations while maintaining realism.





\subsection{Sequential and Controllable LiDAR Scene Generation}

\iffalse
\begin{table}[ht]
\scalebox{1}{
\begin{tabular}{l|ccccc} 
\toprule
         &  MMD$^1$$\downarrow$ & JSD$\downarrow$ & FRD$\downarrow$  & FVD$\downarrow$    \\ 
\midrule
LidarDM\textbf{}      & 25.38 &  0.153  & 1604.10  & 28.41  \\
UniScene\textbf{}      & \underline{21.22} &  \underline{0.141}  & \underline{1007.89}  & \underline{20.28}  \\
\rowcolor{color_ours}
\sysname   &  \textbf{2.13}    &  \textbf{0.049}  & \textbf{292.28}    & \textbf{14.86}  \\
\bottomrule
\multicolumn{5}{l}{\footnotesize 1: MMD has been multiplied by \({10^{4}}.\)} \\
\end{tabular}
}
\centering
\caption{Sequential generation results of different methods on the nuScenes val split.}
\label{tab_seq_con}

\end{table}
\fi



\begin{table}[ht]
\scalebox{1}{
\begin{tabular}{l|ccccc} 
\toprule
         &  MMD$^1$$\downarrow$ & JSD$\downarrow$ & FRD$\downarrow$  & FVD$\downarrow$    \\ 
\midrule
LidarDM\textbf{}      & 25.53 &  0.155  & 1800.18  & 28.48  \\
UniScene\textbf{}      & {21.66} &  {0.143}  & {1182.94}  & {21.04}  \\
\rowcolor{color_ours}
\sysname   &  \textbf{2.94}    &  \textbf{0.069}  & \textbf{743.13}    & \textbf{16.96}  \\

\bottomrule
\multicolumn{5}{l}{\footnotesize 1: MMD has been multiplied by \({10^{4}}.\)} \\
\end{tabular}
}
\centering
\caption{Sequential generation results of different methods on the nuScenes val split.}
\label{tab_seq_con}
\end{table}




\begin{figure}[ht]
    \begin{center}
       \includegraphics[width=\linewidth]{figures/seq_compare.pdf}
    \end{center}
       \caption{Sequential LiDAR scenes generated by different methods on the nuScenes val split. Note that \sysname\ generates sequences of 20 frames, of which we present three representative frames for a clear comparison. (Please refer to the supplementary material for more visual comparisons.)}
       \label{fig_seq_generation}
\end{figure}


We compare \sysname\ with the SOTA sequential LiDAR generation methods, LidarDM \cite{zyrianov2024lidardm} and UniScene \cite{li2024uniscene}. UniScene is a joint RGB-LiDAR generation approach, and we evaluate the quality of its generated LiDAR scenes. It is important to note the difference in sequence length capabilities: LidarDM and UniScene generate sequences of 5 and 8 frames respectively, whereas \sysname\ generates 20 frames at a time. The FVD metric is computed using the the first 5 frames of all the generated sequences to ensure a fair comparison.

\cref{tab_seq_con} presents the results of different methods on the nuScenes val split. It can be seen that \sysname\ achieves lower FRD, MMD and JSD scores than LidarDM and UniScene, indicating improved spatial fidelity and diversity in the generated scenes. In addition, \sysname\ attains the lowest FVD \cite{li2024uniscene} score across the comparing methods, indicating superior temporal consistency in the generated sequences .

\cref{fig_seq_generation} illustrates the qualitative comparison of the sequential LiDAR scene generated by different methods on the nuScenes val split. It can be seen that the objects' structures are mostly broken in UniScene's generated scenes, and the objects' locations diverges from the GT locations  (as evident in the zoomed-in views). In contrast, \sysname\ yields better objects structures and accurate locations, benefiting from the rich information provided by multimodal conditioning. Furthermore, \sysname\ exhibits better temporal consistency with coherent vehicle representations across frames, which is attributed to the EST-Conv and EST-Trans components of \modelname.



\subsection{The Impact of Model Configuration}

We study the impact of EST-Conv and EST-Trans of \modelname\ by evaluating the following model configurations: \textit{Config A} refers to the model where we remove EST-Trans in \sysname\ and replace EST-Conv with regular 2D Convolution blocks. \textit{Config B} refers to the model where we remove EST-Trans in \sysname. \textit{Config C} is the \sysname\ model.

By comparing \textit{Config A} and \textit{Config B} in \cref{tab_ablation_multi}, it is evident that the EST-Conv is critical in enhancing the temporal consistency of the generated LiDAR scenes. By capturing and correlating local features across consecutive frames, the EST-Conv module enables the model to generate temporally coherent and high-quality point clouds. Consequently, both FRD and FVD scores are reduced.
Furthermore, \textit{Config C} yields additional FRD and FVD improvements compared to \textit{Config B}, indicating that EST-Trans substantially enhances spatio-temporal consistency in the generated point clouds. 


\begin{table}
\centering
\begin{tabular}{lcc|cc} 
\toprule
  & \begin{tabular}[c]{@{}c@{}}EST-Conv\end{tabular} & \begin{tabular}[c]{@{}c@{}}EST-Trans\end{tabular} & FRD$\downarrow$ & FVD$\downarrow$ \\ 
\midrule
\textit{Config A} &                                                       &                                                        & 318.05 & 29.50  \\
\textit{Config B} & \checkmark                                             &                                                       & 311.38  & 17.28 \\
\textit{Config C} & \checkmark                                             & \checkmark                                            & \textbf{292.28}   & \textbf{14.86}\\
\bottomrule

\end{tabular}
\caption{
Sequential generation results of \sysname\ with different module configurations.}
\label{tab_ablation_multi}
\end{table}


\subsection{Real-world 3D Object Detection Evaluation}

\begin{table}[h]
\centering
\scalebox{0.98}{
\begin{tabular}{l|cc} 
\toprule
         & mAP $\uparrow$          & GT Agg.$^1$$\uparrow$                  \\ 
\midrule
GT       &   0.804                   & 100.0\%                      \\
UniScene \cite{li2024uniscene} & 0.078 &   9.7\% \\
LidarDM \cite{zyrianov2024lidardm} & 0.140 &   17.4\% \\
\rowcolor{color_ours}
\sysname\ w.o. object priors & 0.123 &  15.3\%  \\
\rowcolor{color_ours}
\sysname\  &\textbf{0.407} & \textbf{50.6}\%  \\
\bottomrule
\multicolumn{3}{l}{\footnotesize 1: GT Agg. denotes aggrements with mAP of GT.} \\
\end{tabular}
}
\caption{3D object detection results on LiDAR scenes generated by different method on the nuScenes val split. (Please refer to the supplementary material for a visualization of detected objects.)}
\label{tab_object_detection}
\end{table}


The accuracy of object localization and the fidelity of objects' points are essential in real-world autonomous driving scenarios, as these factors significantly impact the performance of perception systems.
\cref{tab_object_detection} presents the 3D object detection results for \textit{cars} on LiDAR scenes generated by different methods on the nuScenes val split. It can be observed that UniScene and LidarDM only achieved a mAP of 0.078 and 0.140, respectively, lagging significantly from the results on the real-world GT data. In comparison, \sysname\ achieves a mAP of 0.407, indicating strongest agreement, 50.6\%, with GT data. The ablated variant, \sysname\ w.o. object priors, yields a mAP of 0.123, showing that object priors plays an important role in enhancing object-level fidelities. This is because object priors provide direct and dense guide on objects' point distributions. 
